\documentclass[10pt,a4paper]{amsart}
\usepackage[margin=19mm]{geometry}
\usepackage{amsmath,amssymb,mathtools}
\usepackage[scr=rsfs,cal=euler]{mathalfa}
\usepackage{xfrac}
\usepackage[unicode,hidelinks,bookmarks=false]{hyperref}
\newcommand{\ie}{i.e.\ }
\newcommand{\cf}{cf.\ }
\newcommand{\resp}{respectively\ }
\newcommand{\Subsection}{\subsection}
\providecommand{\nobreakdash}{\nobreak\hbox{-}}
\setlength{\emergencystretch}{2em}
\multlinegap0pt
\usepackage{lmodern}
\usepackage{mathtools}
\usepackage{booktabs}
\usepackage{tabularx,array}
\usepackage{float}
\usepackage{enumitem}
\usepackage{cleveref}
\allowdisplaybreaks
\numberwithin{equation}{section}
\setlength{\emergencystretch}{2em}
\newtheorem{theorem}{Theorem}[section]
\newtheorem{proposition}[theorem]{Proposition}
\newtheorem{lemma}[theorem]{Lemma}
\newtheorem{corollary}[theorem]{Corollary}
\newtheorem{definition}[theorem]{Definition}
\newtheorem{remark}[theorem]{Remark}
\newcommand{\R}{\mathbb R}
\newcommand{\C}{\mathbb C}
\newcommand{\Q}{\mathbb Q}
\newcommand{\Z}{\mathbb Z}
\newcommand{\T}{\mathbb T}
\newcommand{\cS}{\mathcal S}
\newcommand{\cZ}{\mathcal Z}
\newcommand{\cZtwo}{\mathcal Z_2}
\newcommand{\ip}[2]{\langle #1,#2\rangle}
\newcommand{\norm}[1]{\lVert #1\rVert}
\newcommand{\abs}[1]{\lvert #1\rvert}
\newcommand{\e}{\mathrm e}
\newcommand{\ii}{\mathrm i}
\newcommand{\Frob}{\mathrm F}
\begin{document}
\title[An intrinsically subcritical four-point counterexample]{An intrinsically subcritical four-point counterexample}
\author{Secondary 37D30, 65G30 ęywords HRT conjecture, time--frequency shifts, Weyl operators, vector Zak transform, dominated cocycle, validated numerics e̱gin document e̱gin abstract; Vignon Oussa}
\date{}
\maketitle
\noindent\emph{Source and licence.} An intrinsically subcritical four-point counterexample. Version arXiv:2608.07604v1 (CC BY 4.0). This material is distributed under the Creative Commons Attribution 4.0 International licence, http://creativecommons.org/licenses/by/4.0/, as recorded in the arXiv OAI licence metadata for this exact version and shown on the abstract page https://arxiv.org/abs/2608.07604v1. Full rights notice: licenses/arxiv-2608.07604-CC-BY-4.0.txt.\par\smallskip
\noindent\textit{Editorial scope.} Counted unit: the original \texttt{theorem} environment \texttt{thm:validated-certificate} at source lines 476--539 together with its complete proof at lines 541--694; the delivered document keeps source lines 312--871 so that every referenced label and every citation resolves inside it. Native editable source is the paper's own concatenated TeX; the AMS wrapper is editorial formatting only. No publisher class, upstream script or unrelated artwork is redistributed.
\par\smallskip\noindent\textit{Reference note.} This article ships no compiled bibliography (biblatex); the entries delivered here are composed from the author's own BibTeX (.bib) fields, with no field invented and no entry dropped.
\section{Exact vector-Zak reduction}
\label{sec:zak}

For $f\in L^2(\R)$, use the scalar Zak transform
\begin{equation}\label{eq:scalar-zak}
 (\cZ f)(x,\omega)=\sum_{k\in\Z}f(x-k)\e^{2\pi\ii k\omega}
\end{equation}
in the usual $L^2$ sense, and define
\begin{equation}\label{eq:vector-zak}
 (\cZtwo f)_{r+1}(x,\omega)
 =2^{-1/2}(\cZ f)\left(x,\frac{\omega+r}{2}\right),
 \qquad r\in\{0,1\}.
\end{equation}
As in \cite[\S3]{FPPV2026}, this is a unitary map from $L^2(\R)$ onto the
space of measurable $F:\R^2\to\C^2$ satisfying
\begin{align}
 F(x+1,\omega)&=U_1(\omega)F(x,\omega),
 &U_1(\omega)&=\e^{\pi\ii\omega}
                   \begin{pmatrix}1&0\\0&-1\end{pmatrix},
 \label{eq:sewing-x}\\
 F(x,\omega+1)&=U_2F(x,\omega),
 &U_2&=\begin{pmatrix}0&1\\1&0\end{pmatrix}.
 \label{eq:sewing-w}
\end{align}
See also \cite{ZibZee97}.  We call such an $F$ a vector-Zak section.

Put
\[
 D=\begin{pmatrix}1&0\\0&-1\end{pmatrix},\qquad
 S=\begin{pmatrix}0&1\\1&0\end{pmatrix}.
\]
For $m,n\in\Z$, direct use of \eqref{eq:scalar-zak} gives
\begin{equation}\label{eq:half-lattice-action}
 \cZtwo\rho(m,n/2)\cZtwo^{-1}
 =L_{m,n}(x,\omega),\quad
 L_{m,n}=\e^{\pi\ii(nx-m\omega+mn/2)}L_{m,n}^{(0)},
\end{equation}
where
\[
 L_{m,n}^{(0)}=
 \begin{cases}
  \operatorname{diag}(1,\e^{-\pi\ii m}),&n\equiv0\pmod2,\\[1mm]
  \begin{pmatrix}0&1\\ \e^{-\pi\ii m}&0\end{pmatrix},
      &n\equiv1\pmod2.
 \end{cases}
\]
In particular, $L_{1,0}=\e^{-\pi\ii\omega}D$ and
$L_{0,1}=\e^{\pi\ii x}S$.

Let
\begin{equation}\label{eq:base-map}
 \tau=(\alpha,\beta),\qquad Tz=z-\tau,\qquad z=(x,\omega).
\end{equation}
The irrational Weyl shift acts by
\begin{equation}\label{eq:irrational-action}
 (\cZtwo\rho(\zeta)\cZtwo^{-1}F)(x,\omega)
 =\eta(x)F(T(x,\omega)),\qquad
 \eta(x)=\e^{\pi\ii\beta(x-\alpha/2)}.
\end{equation}
Consequently,
\begin{equation}\label{eq:zak-cocycle-action}
 \cZtwo\left[I+\frac35\rho(1,0)+\frac35\rho(0,1/2)\right]
 \rho(\zeta)\cZtwo^{-1}F(z)=B(z)F(Tz),
\end{equation}
where
\begin{align}
 A(x,\omega)&=I+\frac35\e^{-\pi\ii\omega}D
                  +\frac35\e^{\pi\ii x}S,\label{eq:A-definition}\\
 B(x,\omega)&=\eta(x)A(x,\omega).\label{eq:B-definition}
\end{align}
The scalar $\eta$ may be omitted in projective estimates, but not in exact
sewing or multiplier equations.  Direct substitution gives the fiber
covariance
\begin{equation}\label{eq:B-covariance}
 B(z+e_j)=U_j(z)B(z)U_j(Tz)^{-1},\qquad j=1,2,
\end{equation}
where $U_1(z)=U_1(\omega)$ and $U_2(z)=U_2$.

\begin{lemma}[Uniform conditioning]\label{lem:conditioning}
For every $(x,\omega)\in\R^2$,
\begin{equation}\label{eq:conditioning}
 \abs{\det A(x,\omega)}\ge\frac7{25},\qquad
 \norm{A(x,\omega)}_{\mathrm{op}}\le\frac{11}{5},\qquad
 \sigma_{\min}(A(x,\omega))\ge\frac7{55}.
\end{equation}
\end{lemma}

\begin{proof}
One has
\[
 \det A=1-\frac9{25}\e^{-2\pi\ii\omega}
          -\frac9{25}\e^{2\pi\ii x}.
\]
The reverse triangle inequality gives the first estimate, and the triangle
inequality gives the second.  Since
$\sigma_{\min}(A)=|\det A|/\sigma_{\max}(A)$ in dimension two, the third
follows.
\end{proof}

\section{Validated finite-dimensional inequalities}
\label{sec:certificates}

For $N\ge1$, define
\begin{equation}\label{eq:PN}
 P_N(z)=A(z)A(Tz)\cdots A(T^{N-1}z),\qquad
 F_N(z)=\norm{P_N(z)}_{\Frob}^2.
\end{equation}
Let $\sigma_1(M)\ge\sigma_2(M)>0$ be the singular values of an invertible
two-by-two matrix.  Write $\ell_j(M)$ and $r_j(M)$ for its left and right
singular lines.  If $L$ is a complex line and $v\in\C^2$, we use the
phase-independent notation
\[
 \abs{\ip{v}{L}}:=\norm{P_L v},
\]
where $P_L$ is the orthogonal projection onto $L$.  For two lines we similarly
write
\[
 \abs{\ip{L}{L'}}:=\norm{P_L P_{L'}}_{\mathrm{op}},
\]
which is the absolute inner product of any unit representatives.
Throughout, $\ip{u}{v}=u^*v$ is conjugate-linear in the first argument and
linear in the second.

\begin{lemma}[Scalar singular reductions]\label{lem:scalar-reductions}
Let $M\in\mathrm{GL}_2(\C)$, $F=\norm M_{\Frob}^2$, and
$d=\abs{\det M}$.  If $0<\varepsilon<1$, then
\begin{equation}\label{eq:gap-polynomial}
 \varepsilon^2F^2-(1+\varepsilon^2)^2d^2>0
 \quad\Longrightarrow\quad
 \frac{\sigma_2(M)}{\sigma_1(M)}<\varepsilon.
\end{equation}
If $P,Q\in\mathrm{GL}_2(\C)$ have singular ratios $r,s$, if
$0\le\gamma\le1$, and
$X=\abs{\ip{r_1(P)}{\ell_1(Q)}}^2$, then
\begin{equation}\label{eq:junction-identity}
 \frac{\norm{PQ}_{\Frob}^2}
 {\norm P_{\Frob}^2\norm Q_{\Frob}^2}
 =\frac{X(1+r^2s^2)+(1-X)(r^2+s^2)}
 {(1+r^2)(1+s^2)}.
\end{equation}
Thus, if $r,s\le\varepsilon$ and the left side of
\eqref{eq:junction-identity} exceeds
$\gamma^2+2\varepsilon^2+\varepsilon^4$, then $X>\gamma^2$.
\end{lemma}

\begin{proof}
With $r=\sigma_2(M)/\sigma_1(M)$, division of the expression in
\eqref{eq:gap-polynomial} by $\sigma_1(M)^4$ gives
\[
 (\varepsilon^2-r^2)(1-\varepsilon^2r^2),
\]
whose second factor is positive.  For \eqref{eq:junction-identity}, insert
singular-value decompositions and expand the four squared entries.  If
$X\le\gamma^2$, its numerator is at most
$\gamma^2(1+\varepsilon^4)+2\varepsilon^2$, while its denominator is at
least one.  The asserted implication follows.
\end{proof}

Let $R=T^{-1}$ and define the finite stable center
\begin{equation}\label{eq:finite-center-first}
 C(z)=r_2(P_{16}(R^{16}z))=r_2(P_{16}(z+16\tau)).
\end{equation}
The following is the sole computer-assisted input to the analytic proof.


\begin{theorem}[Validated certificate]\label{thm:validated-certificate}
The following statements hold uniformly on $\T^2$.
\begin{enumerate}[label=\textup{(C\arabic*)},leftmargin=2.4em]
\item\label{cert:gap}
For $\varepsilon=10^{-3}$,
\[
 \varepsilon^2F_{16}(z)^2-(1+\varepsilon^2)^2
       \abs{\det P_{16}(z)}^2>3.
\]
\item\label{cert:junction}
For $\gamma=1/4$ and
$c_0=\gamma^2+2\varepsilon^2+\varepsilon^4$,
\[
 F_{32}(z)-c_0F_{16}(z)F_{16}(T^{16}z)>15{,}000{,}000.
\]
\item\label{cert:smooth-overlap}
Let $s_0$ be the smooth step
\[
 s_0(x)=
 \begin{cases}
 0,&x\le0,\\
 \displaystyle\frac{\e^{-1/x}}
 {\e^{-1/x}+\e^{-1/(1-x)}},&0<x<1,\\
 1,&x\ge1,
 \end{cases}
\]
and, as in \cite[\S3.3]{FPPV2026}, let $\chi$ be the smooth vector-Zak
section whose restriction to
$[0,1]\times\R$ is
\begin{equation}\label{eq:smooth-chi}
 \chi(x,\omega)=\frac1{\sqrt2}
 \begin{pmatrix}
  \sin(\pi s_0(x)/2)+\cos(\pi s_0(x)/2)\e^{-\pi\ii\omega}\\
  \sin(\pi s_0(x)/2)-\cos(\pi s_0(x)/2)\e^{-\pi\ii\omega}
 \end{pmatrix}.
\end{equation}
Then
\[
 \abs{\ip{\chi(z)}{C(z)}}>\frac12.
\]
\item\label{cert:linear-gauge}
Let $\chi_{\mathrm{lin}}$ be obtained from \eqref{eq:smooth-chi} by replacing
$s_0(x)$ with $x$ on $[0,1]$ and extending by the sewing laws.  Let $P_C(z)$
be the orthogonal projection onto $C(z)$ and set
\begin{align}
 c_{\mathrm{ref}}(z)&=
 \frac{P_C(z)\chi_{\mathrm{lin}}(z)}
 {\ip{\chi_{\mathrm{lin}}(z)}{P_C(z)\chi_{\mathrm{lin}}(z)}},
 \label{eq:cref}\\
 q_{\mathrm{ref}}(z)&=
 \ip{\chi_{\mathrm{lin}}(z)}
 {B(z)c_{\mathrm{ref}}(Tz)}.
 \label{eq:qref}
\end{align}
Then
\begin{equation}\label{eq:linear-cert}
 \abs{\ip{\chi_{\mathrm{lin}}(z)}{C(z)}}>\frac12,\qquad
 \operatorname{Re}\!\left(\e^{-3\pi\ii/7}q_{\mathrm{ref}}(z)\right)
 >\frac3{20}.
\end{equation}
In \eqref{eq:qref}, the projector at $Tz$ is built from
$P_{16}(R^{16}Tz)=P_{16}(z+15\tau)$.
\end{enumerate}
\end{theorem}

\begin{proof}[Validated verification]
The finite-cover strategy adapts \cite[\S5 and ancillary code]{FPPV2026}.
Every entry of $P_N$ is expanded as a finite Laurent polynomial in
$X=\e^{\pi\ii x}$ and $W=\e^{-\pi\ii\omega}$.  Coefficients, phases, products,
and grid values are evaluated as outward-rounded Arb balls at 160-bit
precision.  For a real trigonometric polynomial
\[
 p(x,\omega)=\sum_{k,l}c_{k,l}\e^{\pi\ii(lx-k\omega)}
\]
put
\begin{equation}\label{eq:fourier-derivative-bounds}
 L_x=\pi\sum_{k,l}|l|\abs{c_{k,l}},\qquad
 L_\omega=\pi\sum_{k,l}|k|\abs{c_{k,l}}.
\end{equation}
On a square cell of side $1/M$ centered at $z_c$,
\begin{equation}\label{eq:cell-cover}
 p(z)\ge p(z_c)-\frac{L_x+L_\omega}{2M}.
\end{equation}
All terms on the right are balls with directed rounding.

For~\ref{cert:gap}--\ref{cert:junction},
\path{arb_fourier_domination_certificate.py} uses a $512^2$ cover.
The grid-center lower bound, derivative allowance, and global lower bound are
\[
\begin{array}{c@{\quad}c@{\quad}c@{\quad}c}
&\text{center}&\text{allowance}&\text{global}\\
\text{gap}&>14.8622&<11.3224&>3.5398\\
\text{junction}&>15{,}955{,}384.58&<555{,}808.77&
>15{,}399{,}575.8.
\end{array}
\]
Residual balls around symbolically cancelling odd modes are charged to both
the value and derivative budgets.

For~\ref{cert:smooth-overlap},
\path{arb_fourier_overlap_certificate.py} proves
\[
 \norm{P_{16}(z+16\tau)\chi(z)}^2
 <\frac{749}{1000}F_{16}(z+16\tau)
\]
on a $1280^2$ cover.  The center lower bound exceeds $591.0561608040093$, the
derivative allowance is $542.179296875$, and the global lower bound is
$48.87686392900937$.  The only non-Fourier input is $0\le s_0'\le2$.
Indeed, after $x=(1-u)/2$,
\[
 s_0'(x)=\frac{2(1+u^2)}{(1-u^2)^2}
 \operatorname{sech}^2\!\left(\frac{2u}{1-u^2}\right)\le2
\]
because $\cosh(y)^2\ge1+y^2$.  Hence
$\norm{\partial_x\chi},\norm{\partial_\omega\chi}\le\pi$.  The coefficient
bounds used in the cell allowance, for
$H=P_{16}(z+16\tau)^*P_{16}(z+16\tau)$ and $F=\operatorname{tr}H$, are
\[
\begin{gathered}
 \norm{\partial_x H}<495000,\qquad
 \norm{\partial_\omega H}<319000,\qquad F<30000,\\
 |\partial_x F|<169000,\qquad |\partial_\omega F|<94000.
\end{gathered}
\]
If $b$ is the squared overlap with the bottom right singular line and
$r=\sigma_2/\sigma_1$, spectral decomposition gives
\[
 \frac{\ip{\chi}{H\chi}}{F}
 =\frac12+\frac{1-r^2}{1+r^2}\left(\frac12-b\right).
\]
Since $(1-r^2)/(1+r^2)\ge\kappa$ from~\ref{cert:gap}, where
$\kappa=(1-10^{-6})/(1+10^{-6})$, the energy inequality
forces $b>1/4$, which is the stated overlap.

For~\ref{cert:linear-gauge}, the fundamental square is divided at
$x=\alpha$ and $\omega=\beta$ into four rectangles.  The exact sewing factors
of $\chi_{\mathrm{lin}}(Tz)$ are fixed in each rectangle.  If
$y=Tz$, $P=P_{16}(y+16\tau)=P_{16}(z+15\tau)$,
$H=P^*P$, $F=\operatorname{tr}H$, and
$\Delta=\sigma_1(P)^2-\sigma_2(P)^2$, then
\begin{equation}\label{eq:Delta-range}
 \kappa F\le\Delta\le F,\qquad
 \kappa=\frac{1-10^{-6}}{1+10^{-6}}.
\end{equation}
The bottom projector is
\[
 P_C=\frac12I+\frac{FI/2-H}{\Delta}.
\]
Put $K_0=FI/2-H$ and
\[
\begin{aligned}
 n_0&=\ip{\chi_{\mathrm{lin}}(z)}
 {B(z)K_0\chi_{\mathrm{lin}}(y)},&
 r_0&=\ip{\chi_{\mathrm{lin}}(y)}
 {K_0\chi_{\mathrm{lin}}(y)},\\
 n_1&=\frac12\ip{\chi_{\mathrm{lin}}(z)}
 {B(z)\chi_{\mathrm{lin}}(y)}.
\end{aligned}
\]
Then
\[
 q_{\mathrm{ref}}(z)=
 \frac{n_0+\Delta n_1}{r_0+\Delta/2}.
\]
The simultaneously certified overlap makes the denominator positive.  With
$\rho_0=\e^{-3\pi\ii/7}$, $\mu=3/20$,
\[
 a=\operatorname{Re}(\rho_0n_0)-\mu r_0,\qquad
 c=\operatorname{Re}(\rho_0n_1)-\mu/2,
\]
the half-plane inequality is equivalent to $a+\Delta c>0$.  It therefore
suffices to check $a+\kappa Fc>0$ and $a+Fc>0$, irrespective of the sign of
$c$.
\path{arb_linear_gauge_certificate.py} uses a $768^2$ cover of each
rectangle.  The least global lower bounds are
$74.94293402046668$ and $74.94319454060898$; the linear-overlap polynomial
\[
 \frac{18}{25}F-
 \ip{\chi_{\mathrm{lin}}(y)}{H\chi_{\mathrm{lin}}(y)}
\]
has global lower bound $792.9113220696389$.
To spell out the overlap deduction, put
$X=\ip{\chi_{\mathrm{lin}}(y)}{H\chi_{\mathrm{lin}}(y)}$ and
$b=\ip{\chi_{\mathrm{lin}}(y)}
{P_C(y)\chi_{\mathrm{lin}}(y)}$.  Since $\chi_{\mathrm{lin}}$ has unit norm,
\[
 b=\frac12+\frac{F/2-X}{\Delta}
 >\frac12-\frac{11}{50\kappa}>\frac14.
\]
Thus $|\ip{\chi_{\mathrm{lin}}(y)}{C(y)}|=\sqrt b>1/2$, as asserted.

On each sewing rectangle, every tested scalar has the exact form
\[
 p(x,\omega)=\sum_{q,r,e}c_{q,r,e}
 \exp\!\left(\pi\ii\bigl[(r/2+e\beta)x-(q/2)\omega\bigr]\right).
\]
Thus
\[
 L_x=\pi\sum_{q,r,e}|r/2+e\beta|\,|c_{q,r,e}|,\qquad
 L_\omega=\pi\sum_{q,r,e}\frac{|q|}{2}|c_{q,r,e}|.
\]
For a rectangle of widths $\Delta x,\Delta\omega$ and an $M^2$ center grid,
the outward-rounded cell allowance is
$L_x\Delta x/(2M)+L_\omega\Delta\omega/(2M)$.  The four chart formulas agree
on their common boundaries by the exact sewing laws.

An independently formulated $1024^2$ implementation, using the opposite
projector-sign convention but sharing the low-level Laurent-product and
Hermitian-polynomial routines, proves the weaker half-plane margin $1/10$.
Its rigorous global numerator and denominator-test lower bounds are
\[
 423.5107387261\qquad\text{and}\qquad2219.0330134840.
\]
This independently checks the source shift, phase, and projector orientation;
it is not an independent software-stack replication or an additional premise.

The theorem states smaller rational bounds than the printed ball endpoints,
so decimal interpretation enters no later deduction.
\end{proof}


By \cref{lem:scalar-reductions,thm:validated-certificate},
\begin{equation}\label{eq:domination-and-junction}
 \frac{\sigma_2(P_{16}(z))}{\sigma_1(P_{16}(z))}<\frac1{1000},\qquad
 \abs{\ip{r_1(P_{16}(z))}{\ell_1(P_{16}(T^{16}z))}}>\frac14.
\end{equation}

\section{The dominated invariant line}
\label{sec:domination}

Define
\begin{equation}\label{eq:BN-KN}
 B_N(z)=B(z)B(Tz)\cdots B(T^{N-1}z),\qquad
 K_N(z)=B_N(R^N z)^{-1}.
\end{equation}
Thus $K_N(z)$ maps the fiber over $R^N z$ to the fiber over $z$.  Since
$B_N$ and $P_N$ differ by a unit scalar, they have the same singular lines
and ratios.  The finite stable center is
\[
 C(z)=\ell_1(K_{16}(z))=r_2(P_{16}(R^{16}z)).
\]
The singular gap makes its spectral projector smooth.  Applying
\eqref{eq:B-covariance} to the full $B_{16}$ product shows projectively that
$C(z+e_j)=U_j(z)C(z)$; no globally phased singular vector is chosen.

\begin{lemma}[Complex projective cone estimate]\label{lem:cone}
Suppose an invertible two-by-two matrix has singular ratio at most
$\varepsilon=10^{-3}$.  Suppose its dominant input line has overlap greater
than $1/4$ with the center of an incoming unitary slope chart.  Then its
projective action maps the disk of radius $\rho=1/100$ in that chart into the
disk of radius $1/200$ about its dominant output line, and has Lipschitz
constant smaller than $1/50$.
\end{lemma}

\begin{proof}
Choose phases so that the incoming center and its orthogonal complement are
\[
 c=\frac{a+qb}{\sqrt{1+|q|^2}},\qquad
 n=\frac{-\overline q\,a+b}{\sqrt{1+|q|^2}},
\]
where $a,b$ are the dominant and subordinate singular input vectors.  The
overlap assumption gives $|q|<\sqrt{15}<4$.  A line represented by $c+wn$
has singular-chart slope
\[
 \Psi(w)=\frac{q+w}{1-\overline q\,w}.
\]
For $|w|\le\rho$, the denominator is nonzero and
\[
 |\Psi(w)|\le\frac{|q|+\rho}{1-|q|\rho},\qquad
 |\Psi'(w)|=\frac{1+|q|^2}{|1-\overline q w|^2}
 \le\frac{1+|q|^2}{(1-|q|\rho)^2}.
\]
In singular coordinates the matrix multiplies slopes by at most
$\varepsilon$.  Therefore the image radius and derivative are bounded by
\[
 \frac1{1000}\frac{4+1/100}{1-4/100}<\frac1{200},\qquad
 \frac1{1000}\frac{1+4^2}{(1-4/100)^2}<\frac1{50}.
\]
\end{proof}

\begin{lemma}[Invariant sections over a translation]\label{lem:section-theorem}
Let $\mathcal D\to\T^2$ be a smooth bundle of closed complex disks with
unitary transition maps, let $f$ be a torus translation, and let
\[
 \Phi_z:\mathcal D_{fz}\longrightarrow\operatorname{int}\mathcal D_z
\]
be a smooth bundle map.  Suppose that, in the unitary disk charts,
\[
 \sup_{z,w}\norm{D_w\Phi_z(w)}\le\kappa<1.
\]
Then $(\Gamma s)(z)=\Phi_z(s(fz))$ has a unique continuous fixed section,
and this section is $C^\infty$.
\end{lemma}

\begin{proof}
The space of continuous sections of $\mathcal D$ is nonempty and complete for
the supremum of the fiber metric.  Unitary transitions make the metric
independent of the chart.  The mean-value inequality makes $\Gamma$ a
$\kappa$-contraction, proving existence and uniqueness.

We justify regularity without differentiating a merely continuous section.
The graph transform acts on $r$-jets.  Once the invariant $(r-1)$-jet is
known, the transformed derivative of exact order $r$ is affine in that
derivative, with linear part
\[
 D_w\Phi_z(s(fz))\circ(Df)^{\otimes r}.
\]
For a translation, $Df=I$ and all higher derivatives of $f$ vanish.  Thus
the order-$r$ jet transform again contracts by at most $\kappa$ and has a
unique invariant continuous jet.

To see that these formal jets are the derivatives of $s$, take coordinate
difference quotients of $s(z)=\Phi_z(s(fz))$.  Translations commute with
difference quotients.  The quotients therefore satisfy affine contraction
equations whose coefficients converge uniformly to the displayed jet
equation.  Stability of fixed points under uniform perturbation proves
convergence to the invariant first jet.  Repeating with higher difference
quotients proves the assertion by induction.  The construction agrees on
chart overlaps.  This is the invariant-section theorem specialized to an
isometric base; compare \cite[Chapter~3]{HPS}.
\end{proof}

\begin{proposition}[Smooth dominated splitting]\label{prop:splitting}
There are smooth line fields $E^s,E^u$ with
$\C^2=E^s(z)\oplus E^u(z)$ such that
\begin{equation}\label{eq:N-invariance}
 B_{16}(z)E^{s/u}(T^{16}z)=E^{s/u}(z).
\end{equation}
The stable line lies within slope $1/200$ of $C(z)$, and the splitting is
uniformly dominated.
\end{proposition}

\begin{proof}
For the stable line, apply \cref{lem:cone} to the inverse graph transform
\[
 (\Gamma_s L)(z)=K_{16}(z)L(R^{16}z).
\]
The needed input junction is the bottom-bottom junction of two consecutive
forward blocks.  More explicitly, put $w=R^{32}z$,
$P=P_{16}(w)$, and $Q=P_{16}(T^{16}w)$.  Because each full $B$-product
differs from its $A$-product by a unit scalar, $K_{16}(z)$ is a unit-scalar
multiple of $Q^{-1}$.  Its dominant input line is therefore $\ell_2(Q)$,
while the preceding inverse-block output center is $r_2(P)$.  In dimension
two,
\[
 |\ip{\ell_2(Q)}{r_2(P)}|
 =|\ip{\ell_1(Q)}{r_1(P)}|>\frac14
\]
by \eqref{eq:domination-and-junction}.  The complete metric space of
continuous sewn sections of the radius-$1/100$ cone bundle is mapped into
itself and contracted by $1/50$.  Indeed, covariance
\eqref{eq:B-covariance} and covariance of the center spectral projector show
that the fiber graph maps agree under the unitary transitions $U_1,U_2$.
Thus \cref{lem:section-theorem} applies: its unique fixed point is a smooth
line $E^s$ in the radius-$1/200$ subcone.

For completeness, construct the second line rather than infer it from the
first.  Let
\[
 C_u(z)=\ell_1(P_{16}(z)).
\]
The forward graph transform
\[
 (\Gamma_u L)(z)=B_{16}(z)L(T^{16}z)
\]
acts on the radius-$1/100$ cone bundle about $C_u$.  Its incoming center
$C_u(T^{16}z)$ has overlap greater than $1/4$ with
$r_1(P_{16}(z))$ by the second inequality in
\eqref{eq:domination-and-junction}.  Hence the same cone lemma produces a
unique fixed line $E^u$, within slope $1/200$ of $C_u$.  The same covariance
check and \cref{lem:section-theorem}, now with the base translation $T^{16}$,
show that $E^u$ is smooth and obeys the sewing transitions.

These lines are transverse.  Indeed, at the source of a block $P_{16}(z)$,
$E^s(T^{16}z)$ has slope at most $1/200$ from $r_2(P_{16}(z))$.  Meanwhile
$E^u(T^{16}z)$ has, in the chart about $r_1(P_{16}(z))$, slope at most
\[
 \frac{4+1/200}{1-4/200}<\frac{21}{5}.
\]
In the $r_1$ chart, the stable disk has slope at least $200$, whereas the
unstable line has slope below $21/5$; hence the lines cannot coincide.  For
the domination estimate, a unit vector in the stable line has the form
$(r_2+w r_1)/\sqrt{1+|w|^2}$ with $|w|\le1/200$, while a unit vector in the
unstable line has the form $(r_1+w'r_2)/\sqrt{1+|w'|^2}$ with
$|w'|<21/5$.  Applying the singular-value decomposition gives
\[
 \frac{\norm{B_{16}(z)|_{E^s(T^{16}z)}}}
      {m(B_{16}(z)|_{E^u(T^{16}z)})}
 \le
 \frac{\sqrt{10^{-6}+(1/200)^2}\,
       \sqrt{1+(21/5)^2}}{1}
 <\frac1{40},
\]
where $m$ denotes the conorm and the common factor
$\sigma_1(P_{16}(z))$ has been cancelled.  Iteration proves uniform
domination.
\end{proof}

\begin{thebibliography}{99}
\bibitem[]{FPPV2026} Faulhuber, Markus; Petersen, Philipp; van Velthoven, Jordy Timo; Voigtlaender, Felix. Linear Dependence of Time--Frequency Shifts of a {Schwartz} Function. (2026)
\bibitem[]{HPS} Hirsch, Morris W.; Pugh, Charles C.; Shub, Michael. Invariant Manifolds. 583 (1977)
\bibitem[]{ZibZee97} Zibulski, Meir; Zeevi, Yehoshua Y.. Analysis of multiwindow {Gabor}-type schemes by frame methods. Appl. Comput. Harmon. Anal. 4 (1997) 188--221
\end{thebibliography}
\end{document}
